\textbf{In-Context Learning.} Pattern reasoning by prompting pre-trained LLMs with few-shot input-output examples is driven by in-context learning \cite{radford2019language,brown2020language}. The examples serve as a form of task specification, where the model is expected to complete further instances of the task by predicting what comes next. In-context learning extends the concept of ``task prefixes'' (predefined token sequences, \eg \cite{raffel2020exploring}), but swapped in with examples instead. \citet{brown2020language} observe that it improves (in particular, out-of-distribution generalization) from scaling model size. This is in contrast to scaling models for pre-training + fine-tuning, which has been shown to not necessarily improve OOD generalization on language tasks \cite{hendrycks2020pretrained}. Nonetheless, despite compelling OOD generalization abilities, in-context learning still comes at a cost, as it continues to lag behind in terms of absolute performance on benchmarks compared to task-specific fine-tuning \cite{radford2019language, dinh2022lift}.

\textbf{Explanations of In-Context Learning.} In-context learning is explicitly trained for by packing examples from the same task and dataset into a context buffer that is fed as input to an LLM with an unsupervised autoregressive objective \cite{brown2020language}, sometimes referred to as meta-training. However, it can also emerge implicitly from training on datasets where tokens exhibit a Zipfian distribution \cite{chan2022data} on Transformer architectures, but not necessarily with recurrent architectures (\eg vanilla RNNs) \cite{chan2022data}. Other works have shown that in-context learning with Transformers can learn simple function classes on par with least squares \cite{garg2022can,akyurek2022learning,von2023transformers}, and can generalize to a seemingly unbounded number of tasks (when trained on tasks from the same task family) better than multitask MLPs \cite{kirsch2022general}, with Bayesian interpretations of this phenomenon \cite{xie2021explanation} \cite{wang2023large}.

\textbf{In-Context vs. In-Weights Learning.} In-context learning occurs during inference without gradient updates to the model weights, and can be differentiated from in-weights learning, which relies on information stored in the model weights during LLM training \cite{chan2022transformers} (and can be useful for completion tasks such as ``Abraham Lincoln was born in \underline{\hspace{1em}}''). \citet{chan2022transformers} observes that generalization of in-context learning can be characterized as more ``exemplar-based'' (on the basis of similarity to in-context examples \cite{shepard1963stimulus}), as opposed to generalization of in-weights learning which tends to be more ``rule-based'' (on the basis of minimal features that support category boundaries in the training data \cite{ashby1986varieties}). The vast capabilities of LLMs \cite{ brown2020language,ouyang2022training,chen2021evaluating,chowdhery2022palm,anil2023palm} have been driven by a combination of both forms of learning. In this work, we are particularly interested in in-context learning, and (depending on the task) using the semantic priors of numeric tokens to drive capabilities such as sequence completion (\cref{sec:sequence_completion}) and improvement (\cref{sec:sequence_improvement}).

\textbf{LLMs and Robotics.} LLMs have been applied across several areas in robotics---such as decomposing high-level task descriptions to mid-level plans \cite{ahn2022can, huang2022language, huang2022inner,zeng2022socratic,song2022llm,huang2023grounded}, robot code \cite{liang2022code,wu2023tidybot,singh2023progprompt,wang2023voyager}, and planning domain definition languages \cite{liu2023llm+}. These methods leverage semantic priors stored in LLMs to compose plans or parameterize primitive APIs, but whether LLMs can directly influence control (\eg at the level of trajectories) in a zero-shot manner remains an open problem. We explore how pattern reasoning capabilities of LLMs may drive various control tasks, to extend or optimize low-level sequences. While it is possible to explicitly train models for these capabilities \citep{laskin2022context, xu2022prompting, zhang2023motiongpt, lee2023supervised}, this work focuses on the inherent abilities of LLMs out-of-the-box, which may have implications for the role of language pre-training for building embodied AI systems. %
Related to our work are \cite{dinh2022lift} which studies how LLMs perform on non-language classification and regression tasks; \cite{webb2023emergent} which examines analogical reasoning in various text tasks; and \cite{brooks2022context} which studies how LLMs can represent a rollout policy and world model in-context and then uses Q-learning to drive policy improvement across a collection of toy environments with linguistic representations. Our use of LLMs for sequence improvement can be seen as a simplification of in-context policy iteration that supports learning from demonstrations and in-context RL, driven by the generality of LLMs as pattern machines. \looseness=-1